\newcommand{\Egreedygap}{27.8}
\newcommand{\EoneTPcert}{0.81}
\newcommand{\EoneFullSec}{541}
\newcommand{\EoneFullGap}{0.01}
\newcommand{\EoneOpt}{11,532}
\newcommand{\LegacyReported}{729}
\newcommand{\LegacyTrue}{16,511}
\newcommand{\LegacyFactor}{23}
\newcommand{\EoneTPgap}{0.06}
\newcommand{\EoneTPsec}{9}
\newcommand{\EoneSavingPct}{30}
\newcommand{\AblText}{The inequalities raise the root LP bound by 14\% and 21\% on the two instances, but within 120\,s HiGHS's own cutting planes and branching make the final bounds of both variants comparable, and the full MILP does not converge faster. Their practical value is different: (SLOT) is required for the aggregated solution to be split into individual vehicles, and without it phase 1 can return plans that no fleet can execute. \Cref{tab:ablation} also shows that relaxing container integrality (phase 1) is what makes the problem easy: the relaxed model solves to optimality in seconds, whereas the integer one does not close in 120\,s.}
\newcommand{\ScalText}{On instances with fewer than 4,000 columns (12 instances) the two-phase scheme is the method of choice: mean gap 3.3\% to the lower bound in 31\,s on average, versus 8.8\% for ALNS started from the greedy plan (60\,s) and 35.6\% for the greedy rule, whose gap is inflated by stock it fails to place within capacities and fleet days. ALNS polishing improves the two-phase plans only marginally (by less than 0.01\%) on these sizes. The picture reverses on the largest class (about 8,100 columns): phase~1 does not converge within its 60\,s budget, its incumbent leaves 114--128\,t unshipped and the restricted phase~2 cannot recover; ALNS from the greedy plan is then the best method (mean gap 21.8\% against the relaxation bound, versus 36.6\% for greedy). The two methods are therefore complementary, and seeding phase~2 with ALNS or greedy routes is an obvious improvement. All gaps are measured against relaxation bounds and overstate the true sub-optimality: on the case study the phase-1 bound was 0.8\% below the proven optimum.}
\newcommand{\SensText}{Internalising a carbon price of 500\,\euro/t cuts emissions from 4.85 to 4.45\,t (8\%) at a cost increase of 0.9\% (\cref{fig:sens}). Scaling the monthly volume from 0.5$\times$ to 3$\times$ moves the cost per tonne from 153 to 701\,\euro{} and the share treated in Spain from 67\% to 44\%, as Spanish capacity saturates; at 3$\times$ the volume 30\,t (lead--acid) cannot be placed within contracted capacities at all, which is precisely the situation the strategic model of \cref{sec:strategic} is designed to anticipate.}
\newcommand{\VSSpct}{6.5}
\newcommand{\StratText}{The stochastic design invests more up front than the expected-value design (22.8 vs 8.5\,M\euro{} of discounted depot and plant costs) and recovers it in operation: second-stage costs fall by 3.2, 7.8 and 54.4\,M\euro{} in the pessimistic, neutral and optimistic scenarios. The expected-value design (Madrid + Bilbao depots; only Leon plant in 2035) leaves 28.1\,kt of Li-ion untreated over the horizon if adoption is fast, versus 2.5\,kt for the stochastic design. EVPI is 12.4\% of the expected cost, an upper bound on what better demand forecasts are worth. Across $P(\text{optimistic})\in\{0.1,0.25,0.4\}$ and plant-cost multipliers $\{0.75,1,1.5\}$, the first Spanish Li-ion plant opens between 2032 and 2034, except when the optimistic scenario is unlikely ($p=0.1$) and plants are not cheaper than assumed ($\times$1, $\times$1.5), in which case none is built.}
\newcommand{\FirstMileText}{On the uniform benchmark, CW+LS is within a few percent of the best-known values reported in the learning-to-route literature (about 6.1, 10.4 and 15.7 for $n=20,50,100$ \citep{kool2019,kwon2020}). \todo{Train POMO on a GPU (\texttt{python -m brlo.firstmile.train}) and add its results; the published POMO numbers suggest gaps of about 1\% to optimal on uniform instances.}}
